% swift-syntax-control-flow-errors.tex — a catalogue of control-flow and error syntax,
% line by line: labeled statements (break /
% continue out of several scopes), guard's scope rule, defer order (LIFO, on every
% exit incl. throw — traced through labeled loops), if-let shorthand, for case let x?,
% for-where, while let, repeat, switch on tuples + where + fallthrough, catch patterns,
% try? / try! / rethrows, typed throws and throws(Never), Result bridging, #available /
% #unavailable, if case vs ~=.
% Deliberately NOT repeated (see Related): the error model, Task error propagation and
% LocalizedError (error-handling), the 8 pattern kinds (swift-enums-pattern-matching),
% if/switch expressions and await-in-defer (parameter-packs-and-new-syntax).
% Sources (checked 2026-10-06) — Swift Evolution, numbers and versions against the swift.org
% evolution index: SE-0345 (if let shorthand, 5.7), SE-0413 (typed
% throws, 6.0: `throws` = `throws(any Error)`, non-throwing = `throws(Never)`, `do
% throws(E)`), SE-0290 (#unavailable, 5.6), SE-0230 (try? flattening, 5.0), SE-0493
% (await in defer, 6.4); The Swift Programming Language — Control Flow (labeled
% statements, fallthrough), Statements (defer: "executed in reverse order"; guard else
% must not fall through), Error Handling.
% @labels: area=swift kind=concept platform=apple topic=language discussion=166
% @tags: labeled-statements, defer, guard, if-let-shorthand, for-case, fallthrough, catch-patterns, typed-throws, rethrows, result-type, available, unavailable
\documentclass[8pt]{extarticle}
\usepackage{printup-sheet}

\lstdefinelanguage{SwiftSheet}{
  morekeywords={protocol,class,final,struct,enum,func,var,let,init,case,
    if,else,return,guard,self,nil,true,false,in,throws,rethrows,try,throw,
    do,catch,defer,for,while,repeat,break,continue,fallthrough,switch,default,
    where,as,is,Never,some},
  sensitive=true, morecomment=[l]{//}, morestring=[b]",
  literate={->}{{\hbox{-}\hbox{>}}}2 {==}{{\hbox{=}\hbox{=}}}2 {!=}{{\hbox{!}\hbox{=}}}2
           {??}{{\hbox{?}\hbox{?}}}2 {...}{{\hbox{.}\hbox{.}\hbox{.}}}3
           {..<}{{\hbox{.}\hbox{.}\hbox{<}}}3 {>=}{{\hbox{>}\hbox{=}}}2
           {<=}{{\hbox{<}\hbox{=}}}2 {+=}{{\hbox{+}\hbox{=}}}2 {/=}{{\hbox{/}\hbox{=}}}2
           {&&}{{\hbox{\&}\hbox{\&}}}2 {||}{{\hbox{|}\hbox{|}}}2}
\lstdefinestyle{cat}{language=SwiftSheet, basicstyle=\ttfamily\scriptsize,
  aboveskip=1pt, belowskip=2pt}

\makeatletter
\DeclareRobustCommand\op[1]{\texttt{\@tfor\@c:=#1\do{\hbox{\@c}}}}
\makeatother
\newcommand\grp[1]{\par\vspace{1pt}\noindent{\color{sheetBlue}\bfseries\footnotesize #1}\par}

\tikzset{
  sc/.style={draw, thick, rounded corners=3pt, anchor=north west},
  ev/.style={draw, thick, rounded corners=1.5pt, minimum height=4.6mm, minimum width=7.4mm,
             font=\ttfamily\tiny, inner sep=0.5pt},
  lbl/.style={font=\tiny, text=black!80, inner sep=1pt, align=center},
  ttl/.style={font=\bfseries\scriptsize, anchor=west},
}
\colorlet{cA}{sheetOrange}\colorlet{cB}{sheetBlue}\colorlet{cC}{sheetGreen!70!black}
\colorlet{cG}{black!45}

\begin{document}

\sheettitle{Swift syntax tricks — control flow, scopes \& errors}{swift · folio}

\oneliner{Every \texttt{\{ \}} is a \textbf{scope}, and every way out of it —
end, \texttt{return}, \texttt{break}, \texttt{continue}, \texttt{throw} — runs that scope's
\texttt{defer}s \textbf{last-in first-out}, innermost scope first. A \textbf{label} names a
loop/\texttt{if}/\texttt{do}/\texttt{switch} so one \texttt{break}/\texttt{continue} leaves
several scopes at once; \texttt{guard} binds into the \emph{enclosing} scope and must leave it;
patterns (\texttt{case let x?}, \texttt{where}) turn loops and conditions into filters.}

\vspace{2pt}
\noindent\begin{minipage}[t]{0.41\linewidth}
\vspace{0pt}
\begin{lstlisting}[style=cat]
enum E: Error { case boom }
func run() throws {
  defer { print("A") }           // fn scope
  outer: for i in 0..<2 {
    defer { print("B\(i)") }     // per outer pass
    for j in 0..<3 {
      defer { print("C\(i)\(j)") } // per inner pass
      if j == 1 { continue outer }
      if i == 1 { throw E.boom }
      print("body\(i)\(j)")
    }
  }
}
do { try run() } catch { print("caught") }
\end{lstlisting}
\end{minipage}\hfill
\begin{minipage}[t]{0.57\linewidth}
\vspace{0pt}
\begin{tikzpicture}[sheet]
  \node[ttl] at (0,2.95) {Nested \texttt{defer}s through a labeled loop — what prints};
  % scopes
  \draw[sc, draw=cA, fill=cA!5] (0,2.7) rectangle (2.75,0.05);
  \node[lbl, anchor=north west, text=cA] at (0.05,2.68) {\textbf{run()} \ defer A};
  \draw[sc, draw=cB, fill=cB!6] (0.12,2.35) rectangle (2.63,0.17);
  \node[lbl, anchor=north west, text=cB] at (0.17,2.33) {\textbf{outer:} pass \texttt{i} \ defer B\textit{i}};
  \draw[sc, draw=cC, fill=cC!8] (0.24,2.0) rectangle (2.51,0.29);
  \node[lbl, anchor=north west, text=cC] at (0.29,1.98) {\textbf{inner} pass \texttt{j} \ defer C\textit{ij}};
  \node[lbl, anchor=north west, align=left] at (0.29,1.62) {exits run defers\\innermost first:\\[1pt]
    end of pass $\to$ C\\\texttt{continue outer} $\to$ C, B\\\texttt{throw} $\to$ C, B, A};
  % timeline
  \foreach \t/\c [count=\k from 0] in {body00/cG, C00/cC, C01/cC, B0/cB,
                        C10/cC, B1/cB, A/cA, caught/sheetRed} {
    \node[ev, draw=\c, fill=\c!10] at (3.05+\k*0.8,2.25) {\t};
    \node[lbl, text=sheetGrey] at (3.05+\k*0.8,2.62) {\the\numexpr\k+1\relax};
  }
  \draw[decorate, decoration={brace, amplitude=3pt, mirror}] (2.69,1.97) -- (4.21,1.97)
      node[lbl, midway, below=3pt, align=center]{i=0 j=0:\\body, then\\end of pass};
  \draw[decorate, decoration={brace, amplitude=3pt, mirror}] (4.29,1.97) -- (5.81,1.97)
      node[lbl, midway, below=3pt, align=center]{j=1:\\\texttt{continue}\\\texttt{outer}};
  \draw[decorate, decoration={brace, amplitude=3pt, mirror}] (5.89,1.97) -- (8.21,1.97)
      node[lbl, midway, below=3pt, align=center, text=sheetRed]{i=1 j=0: \texttt{throw}\\unwinds all 3 scopes,\\innermost first};
  \draw[decorate, decoration={brace, amplitude=3pt, mirror}] (8.29,1.97) -- (9.01,1.97)
      node[lbl, midway, below=3pt, align=center]{caller's\\catch};
  \node[lbl, anchor=west, align=left] at (2.95,0.9) {\texttt{C02}, \texttt{body01} never happen: \texttt{continue outer} skips
    the\\rest of the inner loop. The error leaves \texttt{run()} only \emph{after} A.};
  \node[lbl, anchor=west, align=left, text=sheetBrown] at (2.95,0.35) {\texttt{fatalError()} / a crash runs \textbf{no} defers
    (unwinding,\\not finalizers). A defer can't \texttt{return}/\texttt{throw} its way out.};
\end{tikzpicture}
\end{minipage}

\begin{multicols}{2}
\raggedright

\grp{Labels — leave several scopes at once}
\begin{lstlisting}[style=cat]
outer: for row in grid { for v in row where v < 0 { continue outer } }
loop: while true {
  switch next() {
  case .end: break loop   // plain break would only leave the switch
  default: continue }     // continue always targets a loop
}
parse: do { guard ok else { break parse }; step2() }  // do/if: break only
\end{lstlisting}

\grp{guard — bind for the rest of the scope, or leave it}
\begin{lstlisting}[style=cat]
guard let user else { return }            // user usable BELOW, not inside
guard case .loaded(let items) = state else { return }
guard !xs.isEmpty, let f = xs.first, f > 0 else { throw E.boom }
guard let self else { return }            // in a [weak self] closure (5.7+)
\end{lstlisting}
{\scriptsize \texttt{else} must exit: \texttt{return}, \texttt{throw}, \texttt{break},
\texttt{continue}, or a \texttt{Never} call (\texttt{fatalError}) — falling through does not compile.}

\grp{Optional binding \& loop filters}
\begin{lstlisting}[style=cat]
if let name { greet(name) }        // 5.7 SE-0345: = if let name = name
if let a, let b, a < b { }         // comma = &&, left to right, each binds
if case let .some(x) = opt, x > 0 { }  // pattern + Bool in one condition list
if let n = value as? Int { }       // cast and bind in one step
for case let x? in optionals { }   // [Int?] -> only the non-nil, unwrapped
for case .failed(let m) in results { log(m) }
for i in 0..<n where i.isMultiple(of: 2) { }
while let job = stack.popLast() { run(job) }
repeat { n /= 2 } while n > 1      // body runs at least once
for (i, x) in slice.enumerated() { }  // i = OFFSET from 0, not an index
for t in stride(from: 10, to: 0, by: -2) { }  // 10 8 6 4 2 (through: incl.)
\end{lstlisting}

\grp{defer details · \texttt{do} as a scope}
\begin{lstlisting}[style=cat]
func f() -> Int { var x = 1; defer { x = 99 }; return x }  // returns 1
for url in urls { let h = open(url); defer { close(h) } } // EVERY iteration
func g() async { defer { await h.close() } }  // 6.4 SE-0493: async fns only
do { let tmp = heavy(); use(tmp) }  // plain do = a scope: tmp gone at }
\end{lstlisting}

\grp{Availability — runtime vs compile time}
\begin{lstlisting}[style=cat]
if #available(iOS 17, *) { new() } else { old() }  // RUNTIME OS check
guard #available(iOS 17, *) else { return }
if #unavailable(iOS 17) { legacy() }  // 5.6 SE-0290: `!#available` is invalid
#if os(iOS)                           // COMPILE time: code not built elsewhere
\end{lstlisting}
{\scriptsize \texttt{*} = ``every other platform, from its deployment target''. Only
\texttt{,} combines it with other conditions — no \op{||}, no \texttt{!}.}

\columnbreak

\grp{switch — tuples, where, fallthrough}
\begin{lstlisting}[style=cat]
switch (x, y) {
case (0, 0):              origin()
case (let v, 0), (0, let v): onAxis(v)   // alternatives bind the same names
case let (a, b) where a == b: diagonal()  // where = guard AFTER binding
case (1...9, _):          small(); fallthrough  // run next body, NO match check
case (_, 10...):          edge()
default:                  break           // an empty body must say break
}
if case 200..<300 = status { }            // = if 200..<300 ~= status
\end{lstlisting}

\grp{Errors — catch patterns, try forms, typed throws}
\begin{lstlisting}[style=cat]
do { try load() }
catch let e as URLError where e.code == .timedOut { retry() }
catch DecodeError.missing(let key) { report(key) }
catch { log(error) }                   // implicit `error`
let a = try? load()      // T? - nil on error (T?? flattened since 5.0)
let b = try! load()      // traps on error
func map<T>(_ f: (Element) throws -> T) rethrows -> [T] // throws iff f does
func parse() throws(ParseError) -> AST             // 6.0 SE-0413
do throws(ParseError) { try parse() } catch { fix(error) } // error: ParseError
func pure() throws(Never)      // = cannot throw; throws = throws(any Error)
let r = Result { try load() }  // outcome as a value; try r.get() rethrows
\end{lstlisting}

\section{Traps}
\begin{itemize}
  \trap{\texttt{break} in a \texttt{switch} inside a loop leaves the \emph{switch} — label the loop.}
  \trap{\texttt{defer} reads variables at \emph{exit}, not at declaration; a trailing
        \texttt{defer} as the last statement just runs immediately (warning).}
  \trap{\texttt{fallthrough} skips the next pattern check and can't enter a case that binds.}
  \trap{\texttt{try?} erases \emph{why}; \texttt{nil} from a \texttt{T?} function looks the same.}
  \trap{\texttt{enumerated()} on a slice gives offsets, so \texttt{slice[i]} can crash — use
        \texttt{zip(slice.indices, slice)}.}
  \trap{\texttt{guard let x} inside a closure binds in the \emph{closure}; the \texttt{else}
        \texttt{return} leaves the closure, not the outer function.}
\end{itemize}

\section{Remember}
\textbf{``Exits run defers backwards, innermost first; a label picks the scope; \texttt{guard}
must leave it.''}


\end{multicols}

\noindent{\footnotesize\color{sheetGrey}\textit{Related:} error-handling (model, Task errors,
\texttt{LocalizedError}) · swift-enums-pattern-matching (8 pattern kinds) · swift-optionals ·
parameter-packs-and-new-syntax (if/switch expressions, \texttt{await} in \texttt{defer}) ·
swift-syntax-closures-calls · swift-attributes-directives (\texttt{\#if})}

\end{document}
